\section{Discussion}
\textbf{Why count as well as pixels.} The two evaluations diverge: without colour consensus the reconstruction kept 89\% recall and 80\% precision yet count F1 fell from 0.77 to 0.26, and mask union matched our recall but lost 14--21 points of precision at three to five branches. A method tuned on pixel overlap alone could pick a configuration that looks better and counts worse, so counting goals need a count-based check.

\textbf{Comparison with related work.} Hemming \etal~\cite{hemming2014} found that combining viewpoints raised sweet pepper detectability from at most 69\% to 90\%. For our passive sweep at three to five branches, the share of grapes counted rose from 32--64\% to 80--96\% (labelled grapes matched by a detection, out of 25 per capture; the data of Fig.~\ref{fig:visibility}). The measures differ, so this is an analogy, not a ranking. Our baselines support replacing the assumed fruit depth of synthetic aperture imaging~\cite{vaish2006,kurmi2018} with fruit-only alignment. Learned corrections such as GCNet~\cite{quinones2025gcnet} and occlusion GANs~\cite{kierdorf2022} address what we cannot: fruit that no view saw. Every grape was visible in at least three aligned views at three to five branches, but 13 of 25 were visible in at most two at six branches. The approaches are therefore complementary: observation-backed recovery shrinks the hidden fraction before a learned or statistical correction~\cite{millan2018,nuske2014} is applied.

\textbf{Strengths and limitations.} The method needs one RGB camera and no depth sensor. Its reconstruction varied little across reference views (F1 standard deviation 0.04--0.05, against up to 0.18 for single views), and its gains exceeded the registration noise floor. Its limitations are as follows. All conditions share one arrangement of 25 grapes, so our intervals cover which grapes are present, the reference and the registration, but not variation between scenes. Fruit, foliage and truth masks all come from SAM~3, so GrapeSAM adds independence only at the counting stage. GrapeSAM was trained on field images~\cite{tong2026vivid} and overcounted even the leaf-free photograph. One similarity per view treats the fruit as a single plane, which contributes to cluster offsets and rejected views; the perturbation study bounds global registration error only, not this per-cluster offset, which is one reason the count metric is the primary result. Agreement is tested in raw RGB at a fixed threshold, so exposure or white-balance changes between views, or highlights that move with viewpoint, would reject true fruit; the rig held illumination constant. The joint solve, view validation, foliage mask and baseline rule had no effect at five branches, and their value where views fail is not shown. We tested occlusion (3--6 branches), spacing (20--60~mm), the number of views, the choice of reference and truth registration error. We did not test a second arrangement, green fruit, outdoor light, vehicle motion or yield.

\textbf{Future work.} The priority is repetition across arrangements, including green grapes, with held-out labels. Recovery under heavy occlusion was limited by alignment, so per-cluster or depth-aware alignment is the most promising change; our simulation-only depth-based variant (Sec.~\ref{sec:evolution}) suggests it is feasible but sensitive to pose error. A vehicle test in a vineyard row, with calibration against harvested mass, should follow.
